\subsection{Módulos planos}
\label{subsection:0.6.1-fr}

\begin{env}[6.1.1]
\label{0.6.1.1-fr}
Sea $M$ un $A$-módulo. Las propiedades siguientes son equivalentes:
\begin{enumerate}
  \item[$a)$] El funtor $M\otimes_A N$ en $N$ es exacto en la categoría
    de los $A$-módulos;
  \item[$b)$] $\operatorname{Tor}_i^A(M,N)=0$ para todo $i>0$ y todo
    $A$-módulo $N$;
  \item[$c)$] $\operatorname{Tor}_1^A(M,N)=0$ para todo $A$-módulo $N$.
\end{enumerate}

Cuando $M$ verifica estas condiciones, se dice que $M$ es un
$A$-\emph{módulo plano}. Es claro que todo $A$-módulo libre es plano.

Para que $M$ sea un $A$-módulo plano, basta que para todo ideal
$\mathfrak{J}$ de $A$, de \emph{tipo finito}, la aplicación canónica
$M\otimes_A\mathfrak{J}\to M\otimes_A A=M$ sea \emph{inyectiva}.
\end{env}

\begin{env}[6.1.2]
\EGAReaderTarget{canon-ref.0.6.1.2-fr}{0.6.1.2-fr}
Todo límite inductivo de $A$-módulos planos es un $A$-módulo plano. Para
que una suma directa $\bigoplus_{\lambda\in L}M_\lambda$ de $A$-módulos sea
un $A$-módulo plano, es necesario y suficiente que cada uno de los $A$-módulos
$M_\lambda$ sea plano. En particular, todo $A$-módulo proyectivo es plano.

Sea
\[
  0\longrightarrow M'\longrightarrow M\longrightarrow M''\longrightarrow 0
\]
una sucesión exacta de $A$-módulos, tal que $M''$ es \emph{plano}. Entonces, para
todo $A$-módulo $N$, la sucesión
\[
  0\longrightarrow M'\otimes N\longrightarrow M\otimes N
  \longrightarrow M''\otimes N\longrightarrow 0
\]
es exacta. Además, para que $M$ sea plano, es entonces necesario y suficiente que
$M'$ lo sea (pero puede ocurrir que $M$ y $M'$ sean planos sin que
$M''=M/M'$ lo sea).
\end{env}

\begin{env}[6.1.3]
\label{0.6.1.3-fr}
Sean $M$ un $A$-módulo plano, $N$ un $A$-módulo cualquiera; para dos
submódulos $N'$, $N''$ de $N$, se tiene entonces
\[
  \operatorname{Im}(M\otimes(N'+N''))
  =\operatorname{Im}(M\otimes N')+\operatorname{Im}(M\otimes N'')
\]
\[
  \operatorname{Im}(M\otimes(N'\cap N''))
  =\operatorname{Im}(M\otimes N')\cap\operatorname{Im}(M\otimes N'')
\]
(imágenes tomadas en $M\otimes N$).
\end{env}

\begin{env}[6.1.4]
\label{0.6.1.4-fr}
Sean $M$, $N$ dos $A$-módulos, $M'$ (resp. $N'$) un submódulo de $M$
(resp. $N$), y supongamos que uno de los módulos $M/M'$, $N/N'$ es plano.
Entonces se tiene
\[
  \operatorname{Im}(M'\otimes N')
  =\operatorname{Im}(M'\otimes N)\cap\operatorname{Im}(M\otimes N')
\]
(imágenes en $M\otimes N$). En particular, si $\mathfrak{J}$ es un ideal
de $A$ y si $M/M'$ es plano, se tiene
$\mathfrak{J}M'=M'\cap\mathfrak{J}M$.
\end{env}

